\section{Problem}
\label{sec:problem}

% Your notes, kept as you wrote them:
%
% The Main Problem:
% Navilott is a self-driving car that can navigate through a course with obstacles and reach a destination. The car must be able to detect
% and follow typical road rules --- follow lane lines, obey traffic lights and stop signs at intersections, and safely reach the end of a
% course, with a stretch goal being obstacle avoidance.
%
% The Constraints:
% We are constrained with a budget of \$100 for the project, where we must develop a PCB to be integrated with any of the
% car's other hardware components. Power is limited to up to 12V. Furthermore, the car must be able to interface with a human, where a time
% and battery life is displayed on some screen.
%
% The Course Requirements:
% The course will consist of a 6 intersections on a black mat, with 4 streets coming from each intersection. Each inner street block leads to
% another intersection any outer street block leads to the end of a course. 4 traffic lights will be positioned across different junctions on
% the course while stop signs will fill any remaining intersections/streets.

\subsection{Objective}
\TODO{What must the robot do, in one paragraph? Which PRD rules matter
most to vision: right-lane driving, the exact traffic-light rules
(green to enter, yellow only once inside, never on red), the stop sign
at the exit?}
\TODO{Obstacle avoidance: The PRD's executive summary states it as a must. 
Was it formally downgraded? By whom, and when?}
% Source: PRD, Executive Summary; Product Features Requirements

\subsection{Constraints}
\TODO{Which constraints actually shaped your design decisions, and which
were just boxes to tick? Consider: \$100 BOM, the PCB requirement,
12~V limit, size (10 x 11 x 11~cm), local processing only, start button
and 5~s countdown, elapsed time in tenths of a second.}
\TODO{Battery display: the PRD asks the power system to alert
navigation. Why was battery life moved to the display?}
% Source: PRD, General Constraints; Power; Controller; AutoBot Characteristics

\subsection{The Course}
\TODO{Describe the course the robot was built and tested on. Where does
it differ from the PRD's course, and does each difference make the
vision problem easier or harder?}
% PRD: grid of at least nine 4-way intersections; lights at >= 4;
%      ONE stop sign, at the exit; black foam mats, 3/4" white tape,
%      solid borders, dashed center line; bots drive on the right.
% Your notes: six intersections; stop signs at every unlit
%      intersection.  Repo: vision_stack/docs/course.md

\begin{table}[H]
\centering
\caption{\TODO{caption}}
\label{tab:course}
\begin{tabularx}{\textwidth}{@{}lXX@{}}
\toprule
& PRD & Team course \\
\midrule
\TODO{rows} & & \\
\bottomrule
\end{tabularx}
\end{table}

\subsection{Requirements}
\label{sec:requirements}
\TODO{Turn the PRD's prose into requirements that can each be checked
one way. For each: is it straight from the PRD, or derived (something
vision needs so a PRD rule can be met)? What numbers will you hold
yourself to, and why those numbers?}
\TODO{The PRD's validation guidance asks for each object's detection
field (range, angles, minimum size). Its own requirement, or part of
the detection ones?}
% Starting point: vision_stack/docs/requirements.md (P1-P4, D1-D5, R1-R5).
% Note R5 (motor watchdog) is implemented in src/peripherals/drive.py
% even though requirements.md says it isn't.

\begin{table}[H]
\centering
\caption{\TODO{caption}}
\label{tab:reqs}
\small
\begin{tabularx}{\textwidth}{@{}lXl@{}}
\toprule
ID & Requirement & Source \\
\midrule
\TODO{rows} & & \\
\bottomrule
\end{tabularx}
\end{table}
